\QHeading{第5题　用简便方法计算下面各题。}
\par\noindent


\CalcBlock[方法一：简便计算]{1}{\begin{aligned}[t]&\phantom{={}}\color{black}25\times (4+20)\\&\color{answerred}=25\times4+25\times20\\&\color{answerred}=100+500\\&\color{answerred}=600\end{aligned}}
\par
\MethodBlock{方法二：先处理括号或乘除，再拆因数，凑整计算}{把一个因数拆成两个因数，先凑成整百或整千。}{\begin{aligned}[t]&\phantom{={}}25\times (4+20)\\
&=25\times 24\\
&=25\times (4\times 6)\\
&=(25\times 4)\times 6\\
&=100\times 6\\
&=600\end{aligned}}

\MethodBlock{方法三：先处理括号或乘除，再把特殊因数改成整百或整千}{先把25、125等改写成整百或整千除以一个数。}{\begin{aligned}[t]&\phantom{={}}25\times (4+20)\\
&=25\times 24\\
&=(100\div 4)\times 24\\
&=100\times (24\div 4)\\
&=100\times 6\\
&=600\end{aligned}}


\CalcBlock[方法一：简便计算]{2}{\begin{aligned}[t]&\phantom{={}}\color{black}101\times 22\\&\color{answerred}=(100+1)\times22\\&\color{answerred}=100\times22+1\times22\\&\color{answerred}=2200+22\\&\color{answerred}=2222\end{aligned}}
\MethodBlock{方法二：凑整后补偿}{把接近整十、整百的因数凑整，再补回或减去多算的部分。}{\begin{aligned}[t]&\phantom{={}}101\times 22\\
&=101\times (20+2)\\
&=101\times 20+101\times 2\\
&=2020+202\\
&=2222\end{aligned}}

\CalcBlock[方法一：简便计算]{3}{\begin{aligned}[t]&\phantom{={}}\color{black}125\times 32\times 25\\&\color{answerred}=125\times(8\times4)\times25\\&\color{answerred}=(125\times8)\times(4\times25)\\&\color{answerred}=1000\times100\\&\color{answerred}=100000\end{aligned}}
\MethodBlock{方法二：结合因数，先凑整}{用乘法交换律、结合律，先算容易凑整的两个因数。}{\begin{aligned}[t]&\phantom{={}}125\times 32\times 25\\
&=(125\times 32)\times 25\\
&=4000\times 25\\
&=100000\end{aligned}}

\CalcBlock[方法一：简便计算]{4}{\begin{aligned}[t]&\phantom{={}}\color{black}78\times 12-28\times 12\\&\color{answerred}=(78-28)\times12\\&\color{answerred}=50\times12\\&\color{answerred}=600\end{aligned}}
\MethodBlock{方法二：按运算顺序直接计算（对照）}{本法只作结果核对；本题要求简便计算，前面的简便解法仍保留。}{\begin{aligned}[t]&\phantom{={}}78\times 12-28\times 12\\
&=936-336\\
&=600\end{aligned}}

\CalcBlock[方法一：简便计算]{5}{\begin{aligned}[t]&\phantom{={}}\color{black}76\times 8+24\times 8\\&\color{answerred}=(76+24)\times8\\&\color{answerred}=100\times8\\&\color{answerred}=800\end{aligned}}
\MethodBlock{方法二：按运算顺序直接计算（对照）}{本法只作结果核对；本题要求简便计算，前面的简便解法仍保留。}{\begin{aligned}[t]&\phantom{={}}76\times 8+24\times 8\\
&=608+192\\
&=800\end{aligned}}

\CalcBlock[方法一：简便计算]{6}{\begin{aligned}[t]&\phantom{={}}\color{black}35\times 98\\&\color{answerred}=35\times(100-2)\\&\color{answerred}=35\times100-35\times2\\&\color{answerred}=3500-70\\&\color{answerred}=3430\end{aligned}}
\MethodBlock{方法二：按数位拆分计算}{把一个因数拆成整十、整百数与余下部分之和。}{\begin{aligned}[t]&\phantom{={}}35\times 98\\
&=35\times (90+8)\\
&=35\times 90+35\times 8\\
&=3150+280\\
&=3430\end{aligned}}

\CalcBlock[方法一：简便计算]{7}{\begin{aligned}[t]&\phantom{={}}\color{black}300\div 25\div 4\\&\color{answerred}=300\div(25\times4)\\&\color{answerred}=300\div100\\&\color{answerred}=3\end{aligned}}
\MethodBlock{方法二：先处理括号或乘除，再拆分除数，连续相除}{}{\begin{aligned}[t]&\phantom{={}}300\div 25\div 4\\
&=12\div 4\\
&=12\div (2\times 2)\\
&=12\div 2\div 2\\
&=6\div 2\\
&=3\end{aligned}}

\MethodBlock{方法三：先处理括号或乘除，再利用商不变的性质}{被除数和除数同时乘同一个非零数，商不变。}{\begin{aligned}[t]&\phantom{={}}300\div 25\div 4\\
&=12\div 4\\
&=(12\times 5)\div (4\times 5)\\
&=60\div 20\\
&=3\end{aligned}}

\CalcBlock[方法一：简便计算]{8}{\begin{aligned}[t]&\phantom{={}}\color{black}35\times 101-35\\&\color{answerred}=35\times(101-1)\\&\color{answerred}=35\times100\\&\color{answerred}=3500\end{aligned}}
\MethodBlock{方法二：按运算顺序直接计算（对照）}{本法只作结果核对；本题要求简便计算，前面的简便解法仍保留。}{\begin{aligned}[t]&\phantom{={}}35\times 101-35\\
&=3535-35\\
&=3500\end{aligned}}

\CalcBlock[方法一：简便计算]{9}{\begin{aligned}[t]&\phantom{={}}\color{black}360\div (9\times 5)\\&\color{answerred}=360\div9\div5\\&\color{answerred}=40\div5\\&\color{answerred}=8\end{aligned}}
\MethodBlock{方法二：先处理括号或乘除，再拆分除数，连续相除}{}{\begin{aligned}[t]&\phantom{={}}360\div (9\times 5)\\
&=360\div 45\\
&=360\div (5\times 9)\\
&=360\div 5\div 9\\
&=72\div 9\\
&=8\end{aligned}}

\MethodBlock{方法三：先处理括号或乘除，再利用商不变的性质}{被除数和除数同时乘同一个非零数，商不变。}{\begin{aligned}[t]&\phantom{={}}360\div (9\times 5)\\
&=360\div 45\\
&=(360\times 2)\div (45\times 2)\\
&=720\div 90\\
&=8\end{aligned}}

\par\noindent

\QHeading{第6题　解决问题。}

\Problem{① 李叔叔骑车去旅行，第一天骑行了88千米，第二天骑行了104千米，第三天骑行了96千米。李叔叔这三天一共骑行了多少千米?}
\par


\MethodTitle{方法一}
\Answer{\Formula{\begin{aligned}[t]&\phantom{={}}88+104+96\\&=88+(104+96)\\&=88+200\\&=288\text{（千米）}\end{aligned}}}
\Answer{答：三天共骑行288千米。}

\MethodTitle{方法二：按先后顺序逐次相加}
\Answer{先合计前两部分，再加第三部分；这是与凑整法对照的常规算法。}
\Answer{\Formula{\begin{aligned}[t]&\phantom{={}}88+104+96\\
&=192+96\\
&=288\text{（千米）}\end{aligned}}}
\Answer{答：三天一共骑行288千米。}
\Problem{② 李叔叔和王叔叔合作加工一批零件，李叔叔每小时加工49个，王叔叔每小时加工51个，两人共同加工了6小时就完成。这批零件一共有多少个?}

\MethodTitle{方法一}
\Answer{\Formula{\begin{aligned}[t]&\phantom{={}}(49+51)\times6\\&=100\times6\\&=600\text{（个）}\end{aligned}}}
\Answer{答：这批零件共有600个。}

\MethodTitle{方法二：分别计算两部分，再求和}
\Answer{不先合并每份的数量，而是分别求出两部分的总数，再相加。}
\Answer{\Formula{\begin{aligned}[t]&\phantom{={}}49\times 6+51\times 6\\
&=294+306\\
&=600\text{（个）}\end{aligned}}}
\Answer{答：这批零件共有600个。}
